% GENERATED FILE. Edit canonical lessons, then run npm run generate:books.
% canonical-source-hash: fnv1a64:d83f7772f6cc0bdb
% canonical-lessons: AR-C358-hujja, AR-C358-mabda, AR-C358-burhan, AR-C358-athbata, AR-C358-aqnaa

\chapter{Arguments and proof}
\label{ch:ar-arguments-and-proof}
\hlchaptermodality{358}

\begin{chapteropening}
\textbf{By the end of this chapter:} \emph{I can talk about an argument, a principle and a proof, and say that someone proved or convinced.}

Complete the last lesson of chapter 358: I can talk about an argument, a principle and a proof, and say that someone proved or convinced.
\end{chapteropening}

\section[\texorpdfstring{\ar{حجة} (ḥujja)}{حجة (ḥujja)}]{\ar{حجة} (ḥujja) — an argument}

\label{lesson:AR-C358-hujja}



\begin{quote}
Before the new one: say the Arabic for a sentence, then the Arabic for a barber.
\end{quote}

\subsection*{You'll want to know: \ar{حجة}}
\textbf{\ar{حجة}} — \emph{ḥujja} — \textquotedblleft{}an argument\textquotedblright{}.

\begin{quote}
\ar{قدم} \ar{الطالب} \ar{حجة} \ar{قوية} \ar{في} \ar{النقاش}.
\end{quote}

\begin{quote}
\emph{qaddama aṭ-ṭālib ḥujja qawiyya fī an-niqāsh.}
\end{quote}

\begin{quote}
The student put forward a strong argument in the discussion.
\end{quote}

Plural: \ar{حجج}, \emph{ḥujaj}.

\begin{cousinweb}
\textbf{\ar{حجة}} (\emph{ḥujja}) is built on the native Arabic root \emph{ḥ-j-j}, \textquotedblleft{}to make for; to overcome in argument\textquotedblright{} (Lane). A \emph{ḥujja} is what wins a dispute: an argument, or a proof.
\end{cousinweb}

\begin{grammarlens}[title={What you've built}]
One more word for giving reasons and opinions.
\end{grammarlens}

\subsection*{Your turn}
\begin{itemize}
\raggedright
  \item \emph{Say it:} \emph{ḥujja}
  \item \emph{Say it:} \emph{ḥujja}, once more
  \item \emph{Say it:} \emph{qaddama aṭ-ṭālib ḥujja qawiyya fī an-niqāsh.} — The student put forward a strong argument in the discussion.
  \item \emph{Recall:} say the Arabic for a sentence, then the Arabic for a barber, then say \emph{ḥujja} again
\end{itemize}

\begin{tcolorbox}[breakable,colback=teal!4,colframe=teal!35!black,title={Before you move on}]
What does \ar{حجة} mean? (\textquotedblleft{}An argument\textquotedblright{}.) Say it once more.
\end{tcolorbox}
\section[\texorpdfstring{\ar{مبدأ} (mabdaʾ)}{مبدأ (mabdaʾ)}]{\ar{مبدأ} (mabdaʾ) — a principle}

\label{lesson:AR-C358-mabda}



\begin{quote}
Before the new one: say the Arabic for a barber, then the Arabic for an argument.
\end{quote}

\subsection*{You'll want to know: \ar{مبدأ}}
\textbf{\ar{مبدأ}} — \emph{mabdaʾ} — \textquotedblleft{}a principle\textquotedblright{}.

\begin{quote}
\ar{الصدق} \ar{مبدأ} \ar{مهم} \ar{عند} \ar{أبي}.
\end{quote}

\begin{quote}
\emph{aṣ-ṣidq mabdaʾ muhimm ʿinda abī.}
\end{quote}

\begin{quote}
Honesty is an important principle for my father.
\end{quote}

\begin{cousinweb}
\textbf{\ar{مبدأ}} (\emph{mabdaʾ}) is built on the native Arabic root \emph{b-d-ʾ}, \textquotedblleft{}to begin\textquotedblright{} (Lane). \emph{Mabdaʾ} is a noun of place, a starting point: the idea everything else starts from.
\end{cousinweb}

\begin{grammarlens}[title={What you've built}]
One more word for giving reasons and opinions.
\end{grammarlens}

\subsection*{Your turn}
\begin{itemize}
\raggedright
  \item \emph{Say it:} \emph{mabdaʾ}
  \item \emph{Say it:} \emph{mabdaʾ}, once more
  \item \emph{Say it:} \emph{aṣ-ṣidq mabdaʾ muhimm ʿinda abī.} — Honesty is an important principle for my father.
  \item \emph{Recall:} say the Arabic for a barber, then the Arabic for an argument, then say \emph{mabdaʾ} again
\end{itemize}

\begin{tcolorbox}[breakable,colback=teal!4,colframe=teal!35!black,title={Before you move on}]
What does \ar{مبدأ} mean? (\textquotedblleft{}A principle\textquotedblright{}.) Say it once more.
\end{tcolorbox}
\section[\texorpdfstring{\ar{برهان} (burhān)}{برهان (burhān)}]{\ar{برهان} (burhān) — a proof}

\label{lesson:AR-C358-burhan}



\begin{quote}
Before the new one: say the Arabic for an argument, then the Arabic for a principle.
\end{quote}

\subsection*{You'll want to know: \ar{برهان}}
\textbf{\ar{برهان}} — \emph{burhān} — \textquotedblleft{}a proof\textquotedblright{}.

\begin{quote}
\ar{هذه} \ar{الصورة} \ar{برهان} \ar{على} \ar{كلامه}.
\end{quote}

\begin{quote}
\emph{hādhihi aṣ-ṣūra burhān ʿalā kalāmih.}
\end{quote}

\begin{quote}
This photo is proof of what he said.
\end{quote}

Plural: \ar{براهين}, \emph{barāhīn}.

It takes \ar{على}: \ar{برهان} \ar{على} …, proof of ….

\begin{cousinweb}
\textbf{\ar{برهان}} (\emph{burhān}) is a native Arabic word that Lane files under \emph{b-r-h-n}, \textquotedblleft{}to bring clear proof\textquotedblright{} (Lane). Some early grammarians built it on a three-letter root instead, so the shape of its root is disputed.
\end{cousinweb}

\begin{grammarlens}[title={What you've built}]
One more word for giving reasons and opinions.
\end{grammarlens}

\subsection*{Your turn}
\begin{itemize}
\raggedright
  \item \emph{Say it:} \emph{burhān}
  \item \emph{Say it:} \emph{burhān}, once more
  \item \emph{Say it:} \emph{hādhihi aṣ-ṣūra burhān ʿalā kalāmih.} — This photo is proof of what he said.
  \item \emph{Recall:} say the Arabic for an argument, then the Arabic for a principle, then say \emph{burhān} again
\end{itemize}

\begin{tcolorbox}[breakable,colback=teal!4,colframe=teal!35!black,title={Before you move on}]
What does \ar{برهان} mean? (\textquotedblleft{}A proof\textquotedblright{}.) Say it once more.
\end{tcolorbox}
\section[\texorpdfstring{\ar{أثبت} (athbata)}{أثبت (athbata)}]{\ar{أثبت} (athbata) — to prove}

\label{lesson:AR-C358-athbata}



\begin{quote}
Before the new one: say the Arabic for a principle, then the Arabic for a proof.
\end{quote}

\subsection*{You'll want to know: \ar{أثبت}}
\textbf{\ar{أثبت}} — \emph{athbata} — \textquotedblleft{}to prove\textquotedblright{}.

\begin{quote}
\ar{أثبت} \ar{المحامي} \ar{أن} \ar{الرجل} \ar{كان} \ar{في} \ar{المنزل}.
\end{quote}

\begin{quote}
\emph{athbata al-muḥāmī anna ar-rajul kāna fī al-manzil.}
\end{quote}

\begin{quote}
The lawyer proved that the man had been at home.
\end{quote}

Present: \ar{يثبت}, \emph{yuthbit}.

Followed by \ar{أن}: \ar{أثبت} \ar{أن} …, proved that ….

\begin{cousinweb}
\textbf{\ar{أثبت}} (\emph{athbata}) is built on the native Arabic root \emph{th-b-t}, \textquotedblleft{}to be firm, to stand fast\textquotedblright{} (Lane). It is Form IV, to make firm: to establish a thing, to prove it.
\end{cousinweb}

\begin{grammarlens}[title={What you've built}]
One more word for giving reasons and opinions.
\end{grammarlens}

\subsection*{Your turn}
\begin{itemize}
\raggedright
  \item \emph{Say it:} \emph{athbata}
  \item \emph{Say it:} \emph{athbata}, once more
  \item \emph{Say it:} \emph{athbata al-muḥāmī anna ar-rajul kāna fī al-manzil.} — The lawyer proved that the man had been at home.
  \item \emph{Recall:} say the Arabic for a principle, then the Arabic for a proof, then say \emph{athbata} again
\end{itemize}

\begin{tcolorbox}[breakable,colback=teal!4,colframe=teal!35!black,title={Before you move on}]
What does \ar{أثبت} mean? (\textquotedblleft{}To prove\textquotedblright{}.) Say it once more.
\end{tcolorbox}
\section[\texorpdfstring{\ar{أقنع} (aqnaʿa)}{أقنع (aqnaʿa)}]{\ar{أقنع} (aqnaʿa) — to convince}

\label{lesson:AR-C358-aqnaa}



\begin{quote}
Before the new one: say the Arabic for a proof, then the Arabic for to prove.
\end{quote}

\subsection*{You'll want to know: \ar{أقنع}}
\textbf{\ar{أقنع}} — \emph{aqnaʿa} — \textquotedblleft{}to convince\textquotedblright{}.

\begin{quote}
\ar{أقنعت} \ar{أختي} \ar{أبي} \ar{بشراء} \ar{سيارة} \ar{جديدة}.
\end{quote}

\begin{quote}
\emph{aqnaʿat ukhtī abī bi-shirāʾ sayyāra jadīda.}
\end{quote}

\begin{quote}
My sister convinced my father to buy a new car.
\end{quote}

Present: \ar{يقنع}, \emph{yuqniʿ}.

What you convince someone of comes after \ar{بـ}: \ar{أقنعه} \ar{بـ} ….

\begin{cousinweb}
\textbf{\ar{أقنع}} (\emph{aqnaʿa}) is built on the native Arabic root \emph{q-n-ʿ}, \textquotedblleft{}to be content, satisfied\textquotedblright{} (Lane). It is Form IV, to make content: to bring someone round to an idea.
\end{cousinweb}

\begin{grammarlens}[title={What you've built}]
One more word for giving reasons and opinions.
\end{grammarlens}

\subsection*{Your turn}
\begin{itemize}
\raggedright
  \item \emph{Say it:} \emph{aqnaʿa}
  \item \emph{Say it:} \emph{aqnaʿa}, once more
  \item \emph{Say it:} \emph{aqnaʿat ukhtī abī bi-shirāʾ sayyāra jadīda.} — My sister convinced my father to buy a new car.
  \item \emph{Recall:} say the Arabic for a proof, then the Arabic for to prove, then say \emph{aqnaʿa} again
\end{itemize}

\begin{tcolorbox}[breakable,colback=teal!4,colframe=teal!35!black,title={Before you move on}]
What does \ar{أقنع} mean? (\textquotedblleft{}To convince\textquotedblright{}.) Say it once more.
\end{tcolorbox}
